\chapter*{Conclusion et perspectives}
\markleft{Conclusion et perspectives}
\addcontentsline{toc}{chapter}{Conclusion et perspectives}

Ce stage a permis de traiter un enjeu central d'exploitation : améliorer la pertinence de l'alerting pour renforcer la
réactivité opérationnelle et réduire la fatigue d'astreinte. L'analyse de l'existant a montré que la chaîne technique
en place, bien qu'automatisée, générait un bruit significatif en raison d'alertes souvent orientées causes plutôt que
symptômes.

Le travail réalisé a conduit à formaliser une trajectoire d'amélioration fondée sur des niveaux de criticité,
l'association explicite des alertes aux services impactés et une logique de notification plus ciblée. Cette approche
vise à rendre chaque alerte plus lisible, plus priorisable et plus actionnable pour les équipes en charge de la
production.

Au-delà des aspects techniques, le sujet met en évidence l'importance d'une gouvernance durable de l'alerting :
conventions communes, suivi d'indicateurs de qualité, réévaluation régulière des seuils et capitalisation sur les
retours d'incident.

Les perspectives principales concernent l'affinage des niveaux de criticité selon les services, l'intégration plus
forte avec les runbooks d'exploitation et l'industrialisation du pilotage de la qualité des alertes. À terme,
l'ambition est de disposer d'un système d'alerting fiable et évolutif, réellement aligné avec l'impact utilisateur et
les besoins des équipes d'astreinte.
